(a) By \exref{I.1.1}, the coordinate ring is $k[t]$ or $k[t,t^{-1}]$. The corresponding affine varieties are $\mathbb A^1$ and $\mathbb A^1\setminus\{0\}$, the latter identified with $V(xy-1)$ by $t\mapsto(t,t^{-1})$.

(b) A nonempty proper open subset omits some $a\in k$, so $x-a$ is a nonconstant invertible regular function there. But $k[x]$ has only constant units. The empty subset is plainly not isomorphic to $\mathbb A^1$.

(c) An irreducible projective conic has no singular point: moving such a point to $(0:0:1)$ would make its equation a binary quadratic, which factors over $k$. Choose a point and move it to $(0:0:1)$ with tangent $y=0$. Its equation has the form $yz+ax^2+bxy+cy^2=0$, where $a\ne0$ by irreducibility. Changing $z$ and scaling gives $xy=z^2$. The map $(s:t)\mapsto(s^2:t^2:st)$ has inverses $(x:y:z)\mapsto(x:z)$ on $x\ne0$ and $(z:y)$ on $y\ne0$. Thus the conic is isomorphic to $\mathbb P^1$.

(d) In $\mathbb A^2$ there are disjoint closed irreducible curves, for example $V(x)$ and $V(x-1)$. In $\mathbb P^2$ any two closed irreducible curves meet by \exref{I.3.7}(a). A homeomorphism preserves closed irreducible subsets and their dimensions, so the spaces are not homeomorphic.

(e) If $X$ is affine and isomorphic to a projective variety, then $A(X)=\mathcal O(X)=k$ by (3.4), so $X$ consists of one point.
